% !TEX root = ../main.tex
%% Devoir surveillé n° 4 — chapitres 11 à 13
%% Résultats vérifiés par verifications/eval_ds.py.

\sujetEval{DS4}{MPSI · Devoir surveillé}{Devoir surveillé \numero4}
  {Durée : \textbf{4 heures}\\Calculatrice interdite\\Chapitres 11 à 13}{4}
  {Exercice 1/Développements limités/6 pts,Exercice 2/Suite implicite/4 pts,%
   Exercice 3/Espaces vectoriels/5 pts,Exercice 4/Projecteur et matrices/5 pts}
\consigne{Les quatre exercices sont indépendants. Les développements limités
seront présentés avec leur reste.}

\begin{exoEval}{6 points}{Développements limités}
\begin{questions}
  \item Déterminer le développement limité à l'ordre $4$ en $0$ de
    $\mathrm{e}^{\sin x}$.
  \item En déduire $\displaystyle\lim_{x\to0}\frac{\mathrm{e}^{\sin x}-\mathrm{e}^{x}}{x^{3}}$.
  \item Montrer que $\sqrt{x^{2}+2x}=x+1-\dfrac{1}{2x}+o\Bigl(\dfrac1x\Bigr)$
    en $+\infty$. En déduire l'asymptote de la courbe de
    $x\mapsto\sqrt{x^{2}+2x}$ et la position relative.
  \item Calculer
    $\displaystyle\lim_{x\to0}\Bigl(\frac{1}{\sin^{2}x}-\frac{1}{x^{2}}\Bigr)$.
\end{questions}
\end{exoEval}

\begin{exoEval}{4 points}{Suite implicite}
Pour $n\geqslant2$, on pose $f_n(x)=x^{n}+x-1$.
\begin{questions}
  \item Montrer que $f_n$ admet une unique racine $u_n$ dans $\intOO01$.
    Calculer $u_2$.
  \item Calculer $f_{n+1}(u_n)$ et en déduire que $(u_n)$ est croissante.
    Montrer que $u_n\to1$.
  \item On pose $v_n=1-u_n$. Montrer que $n\ln(1-v_n)=\ln v_n$, puis que
    $n\,v_n\sim-\ln v_n$. En déduire que $v_n\sim\dfrac{\ln n}{n}$.
\end{questions}
\end{exoEval}

\begin{exoEval}{5 points}{Espaces vectoriels}
Soit $E=\R_3[X]$ et $F=\ens{P\in E\mid P(1)=P'(1)=0}$.
\begin{questions}
  \item Montrer que $\varphi:P\mapsto\bigl(P(1),P'(1)\bigr)$ est une
    application linéaire surjective de $E$ sur $\R^{2}$. En déduire que $F$
    est un sous-espace vectoriel de $E$ et donner sa dimension.
  \item Montrer que $F=(X-1)^{2}\,\R_1[X]$ et en donner une base.
  \item Montrer que $G=\R_1[X]$ est un supplémentaire de $F$ dans $E$, et
    décomposer $X^{3}$ selon $E=F\oplus G$.
\end{questions}
\end{exoEval}

\begin{exoEval}{5 points}{Projecteur et matrices}
Dans $\R^{3}$, on note $\mathcal{P}$ le plan d'équation $x+y+z=0$ et
$D=\operatorname{Vect}(e_1)$ avec $e_1=(1,0,0)$.
\begin{questions}
  \item Montrer que $\R^{3}=\mathcal{P}\oplus D$.
  \item Soit $p$ le projecteur sur $\mathcal{P}$ parallèlement à $D$.
    Déterminer $p(x,y,z)$, puis la matrice $M$ de $p$ dans la base
    canonique. Vérifier $M^{2}=M$ et comparer $\operatorname{tr}M$ et
    $\operatorname{rg}M$.
  \item Déterminer une base de $\R^{3}$ dans laquelle la matrice de $p$ est
    $\operatorname{diag}(1,1,0)$, et la matrice de passage correspondante.
  \item Donner la matrice de la symétrie $s$ par rapport à $\mathcal{P}$
    parallèlement à $D$, et vérifier que $s^{2}=\mathrm{id}$.
\end{questions}
\end{exoEval}
